\documentclass[output=paper,colorlinks,citecolor=brown]{langscibook}
\ChapterDOI{10.5281/zenodo.18845596}
\author{Kyriaki Topidi\orcid{}\affiliation{}}
\title[When words destroy worlds]{When words destroy worlds: Online hate speech against religious minorities as harm}
\abstract{This contribution is concerned with the ways online hate speech against religious minorities is conducive to harm towards the groups targeted in both individual and collective terms. It considers in particular how language use can operate as an obstacle to identity formation and a force of oppression against religious minority groups, when used as an instrument of violence. Beginning with an exploration of what online hate speech against religious minorities looks like currently, the analysis will initially focus on a socio-technical analytical framework to interpret online hate speech as harm targeting identity formation and as a persisting factor of vulnerability linked to victimhood. Identity as digital practice combined with the impact of online hate on individuals and entire groups will form the key pillars of analysis. Because of the multi-stakeholder constellation relevant for online content moderation linked to hate speech, the role of technology companies will be then considered within broader regulatory efforts that attempt to map, understand and reverse such harm. The concluding remarks will warn against the use of hate speech towards the creation of social consensus on the acceptability of affronts to the dignity of those with different ethnicity or faith. They will propose instead the reconciliation of cultural difference with structural inequality in the digital space as a strategy against hate speech. The contribution finally will emphasize the need for the deconstruction of histories and discourses against religious minorities that are implicit yet pervasive online in order to reduce the harm(s) incurred. }

\IfFileExists{../localcommands.tex}{
   \addbibresource{../localbibliography.bib}
   \usepackage{tabularx,multicol}
\usepackage{url}
\urlstyle{same}
\usepackage{textcomp}
 
\usepackage{langsci-optional}
\usepackage{langsci-lgr}
\usepackage{langsci-gb4e}
\usepackage{langsci-branding}

% \usepackage{fontspec}
\babelprovide{georgian} % Georgian language
\babelprovide{armenian} % Armenian language
\babelfont[georgian,armenian]{rm}{DejaVu Sans}
\newfontfamily\georgianfont{DejaVu Sans}
\newfontfamily\armenianfont{DejaVu Sans}

\usepackage{dialogue}
\usepackage{attrib}
\usepackage{attrib} % speaker labels in Koca's chapter
\usepackage{phoenician} % alf, bet in Mengozzi's chapter

   

\makeatletter
\let\thetitle\@title
\let\theauthor\@author
\makeatother

\newcommand{\togglepaper}[1][0]{
   \bibliography{../localbibliography}
   \papernote{\scriptsize\normalfont
     \theauthor.
     \titleTemp.
In Saloumeh Gholami \& Ulrich Mehlem (eds.), \textit{Identity formation, language and migration dynamics: Minorities in the MENA region and the German diaspora}, Berlin: Language Science Press [preliminary page numbering]
   }
   \pagenumbering{roman}
   \setcounter{chapter}{#1}
   \addtocounter{chapter}{-1}
}

\newbool{bookcompile}
\booltrue{bookcompile}
\newcommand{\bookorchapter}[2]{\ifbool{bookcompile}{#1}{#2}}


\newcommand{\georgian}[1]{{\georgianfont #1}}
\newcommand{\armenian}[1]{{\armenianfont #1}}

\newcommand{\textarab}[1]{{\arabicfont \beginR #1} \endR}
\newcommand{\texthebrew}[1]{{\hebrewfont \beginR #1} \endR}
\newcommand{\textsyriac}[1]{{\syriacfont \beginR #1} \endR}


\newfontfamily\phoenicianfont{Tuffy Regular.ttf}

   %% hyphenation points for line breaks
%% Normally, automatic hyphenation in LaTeX is very good
%% If a word is mis-hyphenated, add it to this file
%%
%% add information to TeX file before \begin{document} with:
%% %% hyphenation points for line breaks
%% Normally, automatic hyphenation in LaTeX is very good
%% If a word is mis-hyphenated, add it to this file
%%
%% add information to TeX file before \begin{document} with:
%% %% hyphenation points for line breaks
%% Normally, automatic hyphenation in LaTeX is very good
%% If a word is mis-hyphenated, add it to this file
%%
%% add information to TeX file before \begin{document} with:
%% \include{localhyphenation}
\hyphenation{
    par-a-digm
    Ver-to-vec
    Go-go-lin
    sa-gen
    mo-misch
    ant-wor-te-te
    dix-wa-zim
    spre-chen
    B-Verf-GE
    Swe-dish
    re-mains
    schrei-ben
    di-pey-vin
    Chris-ten-sen
    Wes-sen-dorf
    Ja-va-khi-shvi-li
    ana-lysis
}

\hyphenation{
    par-a-digm
    Ver-to-vec
    Go-go-lin
    sa-gen
    mo-misch
    ant-wor-te-te
    dix-wa-zim
    spre-chen
    B-Verf-GE
    Swe-dish
    re-mains
    schrei-ben
    di-pey-vin
    Chris-ten-sen
    Wes-sen-dorf
    Ja-va-khi-shvi-li
    ana-lysis
}

\hyphenation{
    par-a-digm
    Ver-to-vec
    Go-go-lin
    sa-gen
    mo-misch
    ant-wor-te-te
    dix-wa-zim
    spre-chen
    B-Verf-GE
    Swe-dish
    re-mains
    schrei-ben
    di-pey-vin
    Chris-ten-sen
    Wes-sen-dorf
    Ja-va-khi-shvi-li
    ana-lysis
}

   \boolfalse{bookcompile}
   \togglepaper[3]%%chapternumber
}{}

\begin{document}
\maketitle


\section{Introduction}\label{ch.1.2.ss.1}

Hate speech, as a form of public discourse, has been linked to large-scale stigmatization, discrimination and violence against religious \isi{minorities} across the globe \citep{un2023}.\footnote{%
Some of the most prominent historical examples include the Holocaust, the Srebrenica genocide in Bosnia and Herzegovina and the Rohingya crisis in Myanmar.} Hate speech, in these contexts, has often operated as a ``precursor to atrocity crimes, including genocide'' \citep{un2023}.

Online spaces have facilitated the creation and circulation of hateful expressions to the extent even of providing accessible avenues to publicize and disseminate hate crimes in real time \citep{laub2019}. As a result, \isi{hate speech} online has known exponential growth in the last years but still notoriously lacks a universal definition. The Council of Europe's Committee of Ministers Recommendation defines \isi{hate speech} as ``covering all forms of expression which spread, incite, promote or justify racial hatred, xenophobia, anti-Semitism or other forms of hatred based on intolerance'' \citep{coe1997hatespeech}. The difficulty in adopting a widely accepted definition stems in part from the fact that it is complex, if not impossible, to set normative standards about the use of language deemed harmful to both individuals and society more broadly. This may be especially the case when an act of speech incites or amplifies violence by one group against another. The violence connected to such speech is often rooted in internalized beliefs and \isi{attitudes} towards \isi{minority} or minoritized groups. Existing empirical work analyzing, for instance, the use of hashtags online targeting religious \isi{minorities} reveals violent and extremist tendencies (\citealt{aguilera-carnero_azeez2016} on \#jihad). More than that, such usage of language online develops further institutionalized forms of racism and discrimination by demonizing entire communities.

The starting point of the present analysis stems from the growing realization that information and communication, along with the technologies that enhance them, have become essential considerations in democratic deliberations \citep{szakacs_bognar2021} and participation in public affairs as well as integration within diverse societies \citep{uk_law_commission2018}. Addressing two key aspects from a diversity governance  perspective, namely that of the formation of the sociocultural identity of \isi{minority} groups as well as that of their \isi{self-perception} and perception by others, the discussion that follows will consider the implications of the lack of \isi{interculturality} in online spaces. This lack will be illustrated through the growth of online \isi{hate speech} against religious \isi{minority} groups. The analysis will explore how hateful expressions against religious \isi{minorities} within online spaces, understood as a separate category of \isi{hate speech} from a regulatory and normative perspective, within online spaces, create individual and group \isi{harm} and at the same time contribute towards the damaging of social cohesion among groups in diverse societies.  

More specifically, by reversing the frame of \isi{interculturality} within language use, this contribution will consider how violent language use can operate as an obstacle to integration and as a force of oppression against religious \isi{minority} groups. By sharing and spreading expressions of online hate, a dehumanizing process and narratives of unworthiness of citizenship for designated groups and communities are triggered \citep{oboler2013,awan_zempi2020}. Beginning with an exploration of what online \isi{hate speech} against religious \isi{minorities} looks like currently \citep{allen2015,aguilera-carnero_azeez2016}, the analysis will initially focus on a socio-technical analytical framework to interpret online \isi{hate speech}. \isi{Identity} as digital practice combined with the impact of online hate on individuals and entire groups (in the form of \isi{harm}) will form the key pillars of analysis. 

Within this frame, and because of the multi-stakeholder constellation relevant for online content moderation \citep{brown2020}, the role of technology companies will be then considered within broader regulatory efforts that attempt to map, understand and reverse such \isi{harm} \citep{chetty_altathur2018}. Their role within platform governance cannot be underestimated, especially when hate ``campaigns'' become framed as ``identity threats'' within such platforms, provoking emotional reactions from users that platforms and tech companies then use for business purposes \citep{hoffmann2019}. 

Accordingly, the analysis will engage with the most salient features and principles dictating the regulation of the digital ecosystem where ``content is cheap, attention is expensive [and] profit margin is wide'' \citep{grambo2019}. Normatively, it will be argued that the absence of principles such as neutrality, equality/non-discrimination and transparency as crucial components within the socio-technical context of online \isi{hate speech} reinforce a process that allows (if not encourages) the transformation of attacks against ideas to attacks against persons through words, both in their cultural and physical dimensions.

The concluding part of the chapter will argue in favor of the reconciliation of cultural difference with structural inequality in the digital space. It will emphasize the need for the deconstruction of histories and discourses against religious \isi{minorities} that are implicit yet pervasive online. Given how social differences shape the kinds of \isi{harm} experienced, responses to online \isi{harm} regulation need to account for the cultural/social context insofar as they shape processes of digital \isi{belonging} of \isi{minority} groups and their members.

\section{Language as an instrument of violence and oppression in digital spaces}\label{ch.1.2.ss.2}

\subsection{Challenges to the formation of sociocultural identity of minority groups online: The case of religious minorities as victims of hate speech}\label{ch.1.2.ss.2.1}

The internet and the digital space more broadly enhance forms of cultural \isi{belonging} and self-identification, while creating opportunities for new forms of religious \isi{belonging} \citep{campbell2010}. In simpler terms, the online space serves both as a place for creation as well as negotiation of identities \citep[125]{cocq2013}. Digital representation, in this context, can be understood in line with a more basic conceptualization of representation as ``the cultural practices and forms by which human societies interpret and portray the world around them and present themselves to others'' \citep[12]{cloke_crang_goodwin2005}, only applied to the digital space. Following this understanding, digital self-representation creates and reflects reality insofar as it reacts to prevailing power relations \citep[327]{helander2014}. Hate speech disrupts self-identification processes online by causing \isi{harm} to individuals and/or groups victimized due to aspects of their identity.

\largerpage
As importantly, for some religious \isi{minority} groups, such as Muslim groups, digital religious environments can also become spaces of empowerment and learning \citep{bunt2000}. New forms of communication between \isi{minority} community members and renewed opportunities for believing and \isi{belonging} often come, nevertheless, at the cost of exposure to considerable and disproportionate amounts of harmful content online.\footnote{%
See for example \citet[para. 21]{unsr2020}, \citet{faloppa_gambacorta_odekerken2023}, \citet[35 et seq.]{eu_fra2023}. See also \citet{topidi2019_words1}.}
The digital ecosystem can be, at the same time, also used as part of militant jihad by a sub-group of \isi{Muslims} in strategic ways, mainly for propaganda and recruiting purposes \citep[18]{kohler_ebner2019}.\footnote{%
Common methods to do so involve the use of echo chamber effects, the manipulation of algorithms or the hijacking of popular hashtags to increase polarization online and offline. See \citet{davey_ebner2017}; \citet{farwell2014}.} 
In these contexts, online discourses are involved in both reflecting but also constructing social reality. The primary aim in the use of these kinds of discourses becomes to construct idealized perceptions of the in-groups as exposed to existential threat and therefore in need of immediate action.\footnote{%
It should be noted that this is a common process shared with extreme right-wing groups. See \citet[205--207]{muller_harrendorf_mischler2022}; \citet{mischler_muller_geng_harrendorf2019}. The main vehicle to execute these identity management efforts relies on ``ideologies of inequality'' that encourage exaggerated positive linguistic evaluation of the in-group while providing harsh and negative devaluation of the out-groups.} 
These kinds of identity construction processes nevertheless lie outside the focus of the present analysis.

Perceiving the digital space as part of the public sphere creates the possibility of reflecting on the pervasive inequalities among categories of online users, taking particularly into account how one's online presence can amplify messages, opinions or ideas. For religious \isi{minorities}, religion becomes commonly conflated with race and has produced complex counter-reactions in identity politics terms (e.g. research on transnational Muslim feminity as discussed in \citealt{falah_nagel2005}; \citealt{dwyer_shah2009}; \citealt{warren2018}).\footnote{%
Indicative examples from the British context illustrating the conflation of identity markers, in particular religion and ethnic origin are: ``WITH THESE PAKIstanis WOULD FUCK OFF FOR GOOD''; ``Fucking Paki's. Kill them All''; ``They all need the lethal injection dirty vile bastards. It is not just paki muslims, all muslims are scum…Bengali muslims, Indian muslims, \ili{English} muslims etc'' (All examples are drawn from \citealt{awan2016}).} 
Those with visible Muslim identity online are targeted, often repeatedly.  Stigmatization, ``othering'' and stereotyping happen through abusive and provocative language, including offline threats, against the victims and their families. As a result of the effect of such repetitive targeting, the line between the offline and online impact of \isi{hate speech} becomes blurred for them, especially in connection to intimidation and abuse \citep[6]{awan_zempi2015_we_fear}. Hostile comments, racist posts, memes and images, often circulated through fake ID profiles, constitute the basic toolbox for hateful online expressions, particularly around ``trigger events''.\footnote{%
See for example the \#KillAllMuslims twitter hashtag after the Paris shooting in January 2015, targeting the Muslims as a group in reaction to the Charlie Hebdo attacks. The hashtag had been circulated at a low level since 2013 but knew widespread circulation after the 2015 events. The hashtag was later successfully taken over by counter-speakers who moved its focus away from hate (see \citealt{BBCTrending2015}).} 
The use of pictures magnifies the range of abuse and suffering of victims online \citep[20]{awan_zempi2015_virtual_physical}. 

Observing such racist content or being the target of such comments directly invites the inquiry of how online racism is experienced by its victims \citep{nadim2023}. The conflation between ethnic, religious and racial difference by victims that has been already suggested (\citealt[5]{nadim2023}, \citealt{garner_seod2014}) leads to their confusion: In simpler words, it is not always clear to them whether they are targeted on the basis of their religious or ethnic difference. In addition, they experience online hatred on an ethno-cultural basis as inevitable (because it is so widespread), especially for those choosing to make their \isi{minority} identity visible \citep[9]{nadim2023}.

Similarly to the offline world, everyday racism is commonly expressed online through subtle and ambiguous experiences of \isi{stigmatization} \citep[2]{nadim2023}. Micro-aggressions, understood as ``brief, everyday exchanges that send denigrating messages to certain people based on their group membership'' \citep[36]{sue_spanierman2020}, particularly affect racialized individuals. They usually take the form of overt abuse online, everyday othering, dismissive discrimination and victim blaming or claiming reverse discrimination by majority group members.\footnote{%
Characteristic examples of tweets that can be construed as micro-aggressions include: ``@rushanaraali defending @SadiqKhan record as @Mayorof London I can't for the life of me think why. Can you?''; ``CatSmithMP You still forcing Lancashire non-muslim children to eat halal school lunches for votes?''. Both examples are drawn from \citet{harmer_southern2021}.} 
For visibly Muslim users, this signifies that they become cumulatively victimized online in racialized terms (e.g. as ``non-white'') while being also considered as ``threats'' \citep{bravo-lopez2011}.

Still, the correlation between \isi{hate speech} and ethnic \isi{minorities} is neither a novel nor an exclusively digital occurrence: It has been studied since the 1980s. A study by \citet{greenberg_pyszczynski1985} found that perceptions of ethnic \isi{minorities} are negatively affected by contact with \isi{hate speech}. More recently, however, \isi{hate speech} has been linked to the dehumanization of \isi{minority} groups and found to encourage the physical distancing between groups in society \citep{soral_bilewicz_winiewski2018}. Physical distancing inevitably has repercussions for intergroup relations through the development of prejudiced \isi{attitudes}.

The prevailing hateful nature of speech against religious \isi{minority} groups such as \isi{Muslims} online is additionally conducive to the creation of new descriptive norms within online spaces whereby hate against the group(s) in question becomes normalized and as such more impactful \citep[4]{soral_liu_bilewicz2020}. Against the backdrop of the research thesis that the use of digital media has created the conditions for the normalization of violent language about \isi{minorities} \citep{soral_liu_bilewicz2020}, the following section will study the ecosystem of online \isi{hate speech} production and circulation against religious \isi{minorities}. Using the diasporic Muslim groups as a case at hand, the effects of online speech on them will be discussed.

\subsection{The impact of online hate speech on religious minority groups}\label{ch.1.2.ss.2.2}

A robust body of recent literature has shown the breadth of negative representations of these groups in the case of Muslim \isi{minorities} in Europe \citep{awan2016,evolvi2018,evolvi2019,koo2018,moe2019,benigni_joseph_carley2017} the breadth of negative representations of these groups but has emphasized less the causes or the impact of such speech from a diversity governance perspective. The latter, which is the lens within the present contribution, raises the question from a socio-legal point of view of how hateful language online can be linked to offline processes of inclusivity and protection of ethnocultural differences. Religious victimization through online \isi{hate speech} is built around the deliberate rejection of individual personhood contained in hateful expressions.\footnote{%
See, for example, the online ``auctions'' organized against Muslim women in India (cf. \citealt{nabi2022}). More recently, following the terrorist attack of Hamas in Israel on 7 October 2023, the use of the hashtags \#f'''islam, \#stopislam, \#banislam rose by 132\% in their use on X (cf. \citealt{rose_davey2023}).} These expressions produce anxiety, exclusion, and fear, as well as ``shared suffering'' for the groups targeted \citep[143]{walters_paterson_mcdonnell_brown2020}.

Harm is perceived in the context of the present analysis as ``significant emotional distress or physical injury or impairment'' but can also refer to affronts to broader societal norms, values and \isi{attitudes} \citep[201]{powell_scott_henry2020}. When referring to the use of particular forms of language, it is used to depict other individuals as less worthy by virtue of their religious affiliation and has been used to harass and intimidate victims \citep[2]{awan2016}. It affects them by marginalizing and demonizing them. In parallel, such \isi{harm} also defines the experiences of ``\isi{belonging}'' within online spaces in terms of inclusion and exclusion for members of ethnic, racial and religious \isi{minorities} \citep{antonsich2010,yuval-davis2006}.

In more analytical terms, online expressions of religious prejudice against religious \isi{minority} groups entail considerable human cost on the victims \citep[2]{vidgen_yasseri_margetts2021}. For example, in an older 2017 review by Tell MAMA on Islamophobia in the UK, it was found that ``fear of victimisation can limit the activity of British \isi{Muslims} who may avoid using public transport, leaving home, or even leaving their neighbourhoods in which they feel safe'' \citep[11]{tell_mama2017}. Hateful content additionally affects the mental health of victims, their integration and relations with the rest of society due to their systematic stigmatization. 

The impact of \isi{harm} can be cumulative for certain sub-groups within the \isi{minority} groups on the basis of age or \isi{gender}, for instance. Characteristically, women \isi{belonging} to Muslim religious \isi{minorities} are affected more heavily by online \isi{hate speech}. Attacks are extreme and sexualized, creating a silencing effect on them \citep{ballon_vikmane2021}. With user engagement connected to vitriolic comments \citep{wall_street_journal2018}, increase in the usage of platforms leads to increases in targeted advertisement and income for platforms. Such increases, however, come at the cost of digital victimhood, as harmful content has been connected with more use of platforms. Data from the UK, for example, show that Muslim women are systematically victimized to the point where abuse becomes a mundane component of their daily lives \citep[33]{hargreaves2015}. Verbal abuse is characteristically triggered by the wearing of the hijab or the niqab \citep[121]{allen2021}, with those wearing them being attacked even more aggressively both online and offline. The victims' well-being and autonomy are also affected \citep{awan_zempi2015_virtual_physical,awan_zempi2016_affinity}. 

Particularly for Muslim communities, Islamophobic content online is also strikingly harmful for their social and cultural capital: It often extends to the offline world in the form of physical violence, discrimination and marginalization \citep{perrigo2020}. In some cases, such content succeeds in constructing religious narratives that spread fear and anger against religious \isi{minorities} \citep[20]{evolvi2022}. These narratives, while initially alternative, through carefully designed social media strategies become mainstreamed and ultimately ``normalized''.\footnote{%
For an empirical demonstration and additional literature references on the question of the type of impact of \isi{hate speech} on Muslim groups see \citet{schmuck_matthes_paul2017}.} 
Depending on the particular context, there is a general tendency within the most-used platforms, such as Twitter (X) and Facebook (Meta), for right-wing populist parties and outlets (or individuals who identify with them), to use their access to large audiences on social media in order to amplify anti-Muslim \isi{hate speech}.\footnote{%
See the characteristic example of France as discussed in \citet{guerin_fourel2021}.}

Overall, given the noticeable growth of online \isi{hate speech} targeting \isi{Muslims}, by securitizing Muslim communities, online \isi{hate speech} also sets the basis for serious limitations of the freedom of expression, freedom of association and political participation of vulnerable segments of contemporary European societies \citep[21]{ecri2022}. The impact of online \isi{hate speech} at a group-based level has additional implications:  Anti-\isi{minority} religious online \isi{hate speech} reinforces group-based divisions and integration processes \citep[3]{vidgen_yasseri_margetts2021}. In response to targeting and \isi{stigmatization}, experiences of online discrimination can also contribute to the formation of \isi{minority} group defensive self-identification mechanisms that are built around ``oppositional consciousness''. This kind of consciousness can be described as ``an empowering mental state that prepares members of an oppressed group to act to undermine, reform or overthrow a system of human domination'' \citep[4]{mansbridge_morris2001}.

The experiences of such \isi{harm} remain nevertheless largely invisible, despite being raced, gendered but also classed. This is because the legal definitions of \isi{hate speech} frame such expressions with reference to a number of common characteristics such as the content, the tone, the nature of the speech, and its targets, but with less reference to its implications.\footnote{%
For a fuller legal analysis of the prevailing standards on \isi{hate speech} see \citet{topidi2019_words2}.} 
Similarly, within regulatory frameworks, recently also adopted at the European level,\footnote{%
See only indicatively, the EU Digital Services Act 2024.} 
\isi{harm} is conceived as a consequence of individualized and interpersonal violence but in abstraction of its more sociocultural consequences on inclusion \citep{gillett_stardust_burgess2022}. Thus, regulatory solutions may tackle illegal content to some extent but not the legal but harmful types of content.

The efficacy of such legal frameworks is questionable as existing regulatory attempts to reverse \isi{harm} by prioritizing a vision of content moderation are based on ``compensatory individualism'' \citep{walcott2019}. This type of individualism shifts the duty to individual users to regulate their own safety (and visibility) online. In this way, users are asked, for example, to increase the filtering of the content they receive or to undertake identification upgrades to control what they see online. At the same time, platforms and states as co-regulators evade accountability and transparency duties\is{Language!duties} with respect to \isi{hate speech}. The role of tech companies, however, in allowing for the growth of \isi{hate speech} is far from negligeable.


\section{The role of tech companies in spreading hate speech against minority groups}\label{ch.1.2.ss.3}

The growth of religious \isi{hate speech} on social media and major platforms leads to the questioning of how effective policy frameworks and the platforms themselves are in regulating and controlling it. The circulation of harmful content, particularly in the form of \isi{hate speech}, has been tightly linked in the scholarly literature with the development of extremist language \citep[1408]{awan_carter_sutch_lally2023}. The role of anonymity of hate speakers can be singled out, in this respect, as one of the important factors that weakens any sense of responsibility for or perception of the impact of one’s speech on others \citep[1408--1409]{awan_carter_sutch_lally2023}.

For example, narratives developed about \isi{Muslims} online present the group as a security threat, a cultural threat or an economic threat \citep[9]{awan2016}. Visibly Muslim women, in addition, are disproportionately singled out, as mentioned earlier \citep{topidi2024_attacking}. 

The aim of this section is to expose the relationship between the digital ecosystem's governance rules, identity and the social positioning of \isi{minority} groups online in relation to \isi{hate speech} and its impact. Inequality, oppression and social injustice are the underlying axes of analysis of the impact of \isi{hate speech} on religious \isi{minorities}. Under the pretence of neutrality claimed by both platforms (and legal systems) for a long time, it can be argued that \isi{hate speech} against religious \isi{minorities} perpetuates exclusionary policies\is{Language!policies} through the use of online language-engineered violence. In this framework, individuals who either express themselves hatefully or are the targets of such speech should be considered as actors/agents rather than as mere users. Considering them as actors allows us to reveal the vulnerabilities and eventual resilience of those religious \isi{minority} groups targeted by hateful expressions \citep{dale2005}. This labeling is also helpful to illustrate how the digital space does not resolve social stratifications but exacerbates them instead \citep[5]{zheng_walsham2021}.

The concept of Galtung's \textit{culture violence} referring to ``those aspects of culture, the symbolic sphere of our existence that can be used to justify or legitimize direct or structural violence'' \citep[291]{galtung1990} is a useful conceptual frame here: Platforms' moderation approaches to the regulation of \isi{hate speech} against \isi{minorities} such as Muslim groups contribute to ``chang[ing] the moral color of an act from red/wrong to green/right or at least to yellow/acceptable'' \citep[292]{galtung1990} despite their alleged neutrality within online spaces.

More than that, the prevailing model that platforms use at present to respond to the proliferation of \isi{hate speech} relies on standard-setting exercises for acceptable speech online. This kind of standard setting is found within platforms' internal community standards/terms of service. Major platforms mention the protection of certain groups based on attributes, characteristics, or categories when explaining their content moderation approaches.\footnote{%
All relevant information related to X's rules can be found at this link: \url{https://help.twitter.com/en/rules-and-policies/x-rules}. All relevant information related to TikTok's terms of service can be found here: \url{https://www.tiktok.com/legal/page/us/terms-of-service/en}. All relevant information related to Google's terms of service can be found here: \url{https://policies.google.com/terms?hl=en-US.}.} 
While the language of such rules/standards varies across platforms, there is no specific reference to the protection of \isi{minority} groups, or any proactive responses to \isi{harm} that they might experience or have experienced. Rather, all platforms rely on a level of self-reporting and have various rankings on which forms of online violence are worse than others. These platforms, within their terms of service, cover presently, though, the bare minimum of what is required of them by legislative frameworks in force.

Top-down centralized approaches to moderation as framed within terms of service, may reduce \isi{hate speech}, but at the same time, they can still have a negative, silencing impact on historically marginalized groups. Content removal as well as the ban of users/communities online may also have an immediate impact on \isi{harm} reduction. Such measures do not address, however, the disruptions to social cohesion and discrimination experienced by \isi{minority} groups. The lack of transparency in moderation decisions adds a further obstacle to how platforms try to build communities of users based on fairness. Standardized rules of moderation where algorithms (predictably) fail to acknowledge cultural nuance and more sophisticated forms of racism increase experiences of intersectional \isi{harm} among \isi{minority} groups in ways that lead to \textit{automated powerlessness}.

In some cases, state regulation becomes an additional layer of normativity (e.g. within the NetzDG in Germany, the Online Safety Act in the UK or the Digital Services Act at EU level). The degree of success of this accumulation of levels of norms remains contested: While on the one hand state regulation may reduce the intensity of \isi{hate speech} \citep{andres_slivko2023}, on the other hand, platforms have the propensity to keep harmful and extreme content online in order to entertain high levels of user engagement and by extension profits from advertising \citep{liu_yildirim_zhang2021}. This is because hateful expressions have been found to receive higher user engagement \citep{mallipeddi_janakiraman_kumar_gupta2021}.

Attempts to regulate online \isi{hate speech} have highlighted two further important dimensions related to the \isi{harm} caused to those targeted by them: First, there seems to be a strong connection between offline and online hate. Hateful and harmful content grows ``by contagion'' as users get exposed to it and normalize it \citep{beknazar-yuzbashev_et_al2022}. Conversely, lesser amounts of hateful content lead algorithms shaping users' feeds to less inclusion and reproduction of \isi{hate speech}. Second, the design of any system of regulation matters: Treating all platform users uniformly within a given regulatory framework creates the incentive for platforms to punish/moderate the content that has comparatively lower impact while showing more deference to higher-impact content. Rather, the type/degree of \isi{harm} caused would represent a more effective regulatory tool to identify and eliminate \isi{hate speech} \citep[3]{andres_slivko2023}.

When defining \isi{harm}, in connection to \isi{hate speech}, platforms tend to adopt limited and subjective definitions of the concept, as already mentioned. Usually approached as physical violence built around notions of danger, \isi{harm} is represented within platforms' policies\is{Language!policies} as a narrow, bilateral nexus between a hate speaker and an individual victim, limiting how such \isi{harm} can be understood in its core and its manifestations \citep[64]{decook_cotter_kanthawala_foyle2022}. And yet the role of platforms within the process of governance of the digital space extends beyond the moderation of content to that of the behaviour of users online. By accommodating racialized and gendered lenses, vulnerable users and \isi{minority} group members in particular are left without meaningful means of protection online \citep[66]{decook_cotter_kanthawala_foyle2022}.

Fuzzy understandings of \isi{harm} create the space to exclude \isi{minority} groups from protection and allow tactical maneuvers from platforms when there is public outcry. Platforms' understandings of \isi{harm} are equally non-binding and malleable to different situations \citep[70]{decook_cotter_kanthawala_foyle2022}. Likewise, the hierarchization of \isi{harm} according to quantitative criteria (e.g., according to the number of times a harmful post has been viewed) is essentializing the more complex ways in which \isi{minority} group members experience it online. The example of harmful stereotyping online is characteristic: As stereotyping does not necessarily or directly lead to offline violence, it is not considered the focus of content moderation, especially because it is rooted in offline behaviours and social understandings. By removing the cultural context from \isi{harm}, religious \isi{minority} groups are exposed to superficial treatment of their digital positionality, often under the pretext that no causal link can be established between online \isi{hate speech} and offline violence and oppression \citep[72]{decook_cotter_kanthawala_foyle2022}. In other words, the conception of \isi{harm} is not neutral in its digital dimension, especially when content ``moderation can be manipulated to propagate values or cultivate doubt'' \citep[75]{decook_cotter_kanthawala_foyle2022}.

Although online hate varies in its circulation across time, context, geography and users, it is worth noting that often a numerically smaller core of extremely active users is responsible for important amounts of such content where the far-right has been very present (\citealt{kreissel_et_al2018,vidgen_yasseri_margetts2021} for findings in Germany and the UK). This suggests that user-driven approaches to online content moderation should not be excluded, especially in light of the large amounts of content in need of moderation. Community-based moderation models have been additionally explored as an alternative in this direction. They propose to remedy the current uneven distribution of power between platforms and users. By relying on a small and more compact group of individuals to moderate content, these approaches to content moderation nevertheless still presuppose that victims formulate complaints in cases of \isi{hate speech} \citep{hasinoff_schneider2022}.

\section{Digital belonging through contextual approaches to hate speech}\label{ch.1.2.ss.4}

\largerpage
One of the direct implications of digital \isi{harm} and victimhood against \isi{minority} religious groups and their members concerns the interlinked processes of identity building and \isi{belonging}. The online experiences of vulnerable users, such as \isi{minority} religious group members that are systematically victimized due to their visible differences, affect their sense of (digital) \isi{belonging}. \textit{Digital belonging} \citep{marlowe_bartley_collins2017} is connected to one's experience in socio-spatial terms between inclusion and exclusion. Such \isi{belonging} can be critically linked to experiences of \isi{hate speech} insofar as the latter shed light on how the (un)-making of \isi{belonging} is performed in virtual environments \citep[86]{marlowe_bartley_collins2017}. 

While digital users, often labelled as ``digital subjects'', represent ``abstractions'' of themselves as far as tech companies are concerned,\footnote{%
The abstraction is largely due to the treatment of individual users by platforms and big tech as subjects to which a piece of personal data relates, not as cultural entities. By extension, digital infrastructures adopt in their regard totalizing tendencies, where context and historicity only relate to business-related purposes and behavior \citep[2--3]{goriunova2019}.} in reality, there has been a wealth of frameworks and discourses that allow us to explore their identity.\footnote{%
See indicatively only \citet{cover2023}.} 
As identity markers condition \isi{identity performance} within the internet, they have become necessary to study diversity but also inclusion/exclusion online.

\isi{Identity} performance online has layers. The shaping of \isi{belonging} is an inherently intersectional process, related to ethnicity/race, religion, and \isi{gender} but also literacy, economic status or education level, to name some of the factors. In cultural terms, it is multi-layered and is experienced within specific socio-spatial terms between inclusion and exclusion \citep{yuval-davis2006}. The type(s) of identity marker(s) that users put forward as prevailing in each case when expressing themselves online, in circumstances where their identities are multi-layered, depend(s) on the context and content of the digital medium \citep[624]{pennington2018}. These experiences matter because digital inclusion and social cohesion are explicitly interlinked: Inequality and vulnerability within society is today increasingly conditioned by digital interactions \citep{marlowe_bartley_collins2017}. Digital media and platforms, therefore, tend to reflect and, in some cases, amplify existing social inequalities. 

By extension, given how social networks and platforms condition social relationships, identity formation online becomes inseparable from cultural identity (i.e., how one perceives themselves but is also perceived by others) \citep[6]{cover2023}. The articulation of ethno-cultural identities is additionally of a relational kind: This means that it becomes dependent on the acknowledgement of similarity or difference \citep[6]{cover2023}. Still, digitization operates against the backdrop of already existing structures of power, hierarchies and imbalances. The deterring use of language, in the form of \isi{hate speech}, in reaction to digital self-identification processes of \isi{minority} group members comes at the cost of exclusion from technological networks, considered as ``one of the most damaging forms of exclusion'' \citep[3]{castells2001}. 

Responses to digital \isi{harm} issues in regulatory terms, however, are largely technical so far (for example, relying on automated or quasi-automated content moderation tools to prevent/suppress \isi{harm}), ignoring the social implications of such harmful speech \citep{gillett_stardust_burgess2022,topidi2024_attacking}. By privileging the individual responsibility of users (either as hate speakers or as victims of \isi{hate speech}) to remain safe from \isi{harm}, the \textit{datafication}\footnote{%
Datafication is here understood as ``a means to access (…) and monitor people's behavior'' (cf. \citealt{van_dijk2014}: 198).} 
of the circulation of harmful content combined with covert and subjective decision-making on the types of content that can be removed/rendered less visible or searchable online remain ineffective.  Characterized aptly as a ``surface approach'' \citep[8]{gillett_stardust_burgess2022}, efforts of platforms and big tech privilege visibility over social change. Typical examples of this approach are the reactive decisions taken by platforms announcing policy updates, conducting stakeholder consultations or those related to the careful release of audits and reports, and the funding of projects in the face of criticism against their responses to \isi{harm} perpetrated online. 

The focus on the removal of individual content suggests that the criteria of recognition and repair of such \isi{harm} are limited in their perspective and remain largely reactive: Emphasis on numerical thresholds, tiered systems of classification of \isi{hate speech} or imperfect algorithmic detection tools inevitably lead to the hierarchization of the kinds of harms that merit the platforms' reactions and attention. 

Despite the general reduction of \isi{hate speech} in some instances, \isi{minority} groups continue to be impacted disproportionately by centralized approaches to moderation. In practice, \isi{hate speech} as harmful content creates complex forms of distress (both physical and psychological), marginalization and exploitation of its victims, to name a few \citep{schoenebeck_blackwell2021}.  Both individual-level and community-level sanctions against harmful content online affect marginalized groups in multiple ways \citep[23]{gilbert2023}. Therefore, applying content moderation policies\is{Language!policies} without nuance may paradoxically end up eliminating content calling out discrimination, anti-\isi{minority} speech and oppression while allowing hateful speech par excellence to remain online \citep{marshall2021}. In such ways, powerful speakers are able to maintain their position in digital ecosystems, while vulnerable speakers are silenced. With racist speech thriving on social media and the evolution of the tactics towards digital \isi{harm} moving from words to images, emojis and other seemingly neutral practices, it is clear that digital technologies both enable and reshape oppression based on features such a race, religion, \isi{gender} as well as intersectionally \citep[206]{matamoros-fernandez_farkas2021}. Within this frame, \isi{hate speech} detection calls for a contextual exercise, where apart from culture, digital interfaces and algorithmic tools, user behaviour plays a crucial role. Overtly abusive, degrading or dehumanizing expressions are only one part of the nexus between digital spaces and structural injustices.

Still, the solution to online \isi{hate speech} and the \isi{harm} it provokes does not lie, as big tech companies claim, within content moderation. The latter is a costly and imperfect process but still one that sustains the profit-making profile of tech companies. It is less aligned with the principles of accountability, digital equality and restorative processes on both the individual and collective levels. Removing hateful expressions and occasionally banning hate speakers from platforms for violations of their internal rules are symptomatic of a partial and incomplete understanding of the problem \citep{salehi2020}. 

To remedy the \isi{harm} caused by \isi{hate speech}, stakeholders need to account for the multiple axes of domination that prevail within digital spaces. More intersectional approaches to \isi{hate speech} would reveal the multi-layered oppression experienced by Muslim women online, for example, framed within institutional, structural and political power structures offline \citep{bivens2017,bowleg2008,christian_day_diaz_peterson-salahuddin2020,macmillan_cottom2020}. Given how the online spaces reflect the offline ones, intersectionality can be useful to lay out how online \isi{harm} and \isi{hate speech} limit \isi{minority} (digital) self-representations through a variety of possible constellations including but not limited to patriarchy, white supremacy, sexism, racism, religious prejudice or classism \citep[44-45]{topidi_metcalfe2024}. Following Gray's intersectional tech framework, decisions in moderation contexts must account for how religious \isi{minorities}, along with other vulnerable groups, are adversely affected by uniform guidelines or victimized as visible targets due to their differences \citep[6]{gilbert2023}. The \isi{harm} experienced by members of religious \isi{minorities} requires addressing the victims' needs as they relate to the harmful effects of such speech, while involving the offenders and the communities in the process of restoring social cohesion and the dignity of those harmed. To respond to these systemic and structural shortcomings, design and policies\is{Language!policies} surrounding \isi{hate speech} moderation require a careful scaling of transparency, accountability and sensitivity to existing matrixes of domination and power differentials \citep{hill-collins1990}.

Victims of \isi{hate speech} are, ultimately, left only with individual resilience remedies to overcome \isi{harm} caused by exposure to \isi{hate speech}. Resorting to resistance and resilience strategies, however, suggests the denial of the symbiotic relationship between technology and people in its socio-technical dimensions. At the same time, this denial discards the complex and multiple identities of religious \isi{minority} group members, especially in the ways that those identities are performed online. The example of Muslim \isi{minority} women is once more characteristic in this respect: Their faith is used as a targeting mechanism. They are singled out as victims of \isi{hate speech} due to their visible signs of otherness. Against these attacks, they develop hybrid practices and identities online to bridge their experience of discrimination with the imperative of integration within mainstream (i.e. majority) societies \citep[65]{mahmudova_evolvi2021}.  

Ultimately, the platforms' language of inclusion, as \citet{hoffmann2021} argues, has fallen short of providing protection and safe visibility for users. Rather, it is wrapped up in positivist and responsive language that downplays the severity of online \isi{harm} through limited inclusion. Therefore, the inclusion discourse used by platforms further increases vulnerability which leads to experiences of ``data violence''.\footnote{%
An example of violence that has been caused by data can be historically reflected in instances where people have been killed because of government storage of data, for example the use of Belgian colonial registration lists of Hutus and Tutsis during the Rwandan genocide.} 
More recent forms of data violence would include the development of surveillance systems that label and classify people in ways that might reproduce structural and symbolic violence towards marginalized groups. Hoffman complicates data violence further by adding algorithmic violence, coined by Mimi \citet{onuoha2018}, to describe how ``an algorithm or automated decision-making system inflicts violence by preventing people from meeting their basic needs''.\footnote{%
For example, in what Ruha \citet{benjamin2019} terms the New Jim Code, transportation apps have tracked and coded areas that are most populated by people of colour as ``dangerous''.}

\section{Concluding remarks: Drawing the barriers of inclusion in digital spaces and the offline world}\label{ch.1.2.ss.5}

As offline spaces become blended within online ones and vice versa \citep{siuda2021}, they create together hybrid spaces constructed around networks which are both decentralized and flexible (in line with Castells's network theory, \citealt{castells2000}). These networks are highly dependent on the structures and actions of big tech where an important and constantly growing amount of network-construction takes place. At the same time, it is becoming clearer that Islamophobic online discourses impact offline hate crimes \citep{levin2016,muller_schwarz2020}. Drawing the two trends together, it can be argued that both online and offline \isi{hate speech} confirm the internalization of certain cultural discourses propagated, including by some politicians \citep{burnett2017}. Far-right and populist groups' online discourses function as enabling factors towards the creation of digital ecosystems, reinforcing exclusion \citep{perry_scrivens2018}. They do so by promoting narratives of danger and untrustworthiness \citep{uk_home_office2020}. 

Notwithstanding concerning patterns, no strategy, policy or targeted legislative framework that tackles anti-Muslim online \isi{harm} exists. This may be because the majority of Islamophobic hate crime, including in its digital dimensions, is considered ``low level'' (e.g. verbal abuse) \citep[118]{allen2021}.\footnote{%
Far-right, anti-Muslim organizations, such as Britain First, have for example switched their focus on Islamophobia from religion to race in order to ``defend'' statues and memorials connected to colonialist accusations \citep[42]{mulhall_khan-ruf2021}. } 
Still, the correlation between online \isi{hate speech} on race/religion with police-recorded religiously aggravated crimes remains positive \citep{williams_burnap_javed_liu_ozalp2020}, although it is unclear whether online \isi{hate speech} precedes or follows offline hate crime.\footnote{%
Other factors such as demographic factors, unemployment rates, etc matter as well \citep[112]{williams_burnap_javed_liu_ozalp2020}.}

Understanding the effects of \isi{hate speech} within digital platforms in the absence of comparable data on their impact on individuals that use them remains a challenge (\citealt[17]{belfer_shorenstein2023}, \citealt[8]{vidgen_margetts_harris2019}). Causality between \isi{hate speech} and hate crimes (i.e. the establishment of a clear relationship of cause between \isi{hate speech} and specific harms) is also largely missing, especially as online experiences may vary across individuals and the interplay between the offline and the online worlds is blurred \citep[20]{belfer_shorenstein2023}. More clarity on the relationship between \isi{hate speech} and crimes could enhance the development of social consensus on why reversing \isi{hate speech} both online and offline represents a worthwhile public policy objective. It could also mitigate the withdrawal and self-censorship of targeted victims from the digital public space.\footnote{%
See for example in the \ili{German} context the findings of the GMK (2024)\nocite{gmk2024}, \url{https://kompetenznetzwerk-hass-im-netz.de/lauter-hass-leiser-rueckzug/}.}
This kind of clarity is particularly relevant  for a number of contexts where \isi{hate speech} online has been found to be produced and distributed by a limited number of users with the intention to create a distorted representation of public opinion.\footnote{%
Indicatively, in a 2018 report, 5 per cent of hateful speech accounts generated 50 per cent of likes for degrading comments in Germany (cf. \citealt{kreissel_et_al2018}).}

By neglecting the needs and experiences of the targets of \isi{hate speech}, but also the socio-technical implications of such kinds of victimhood based on religious identity markets (though not only), inclusion in the digital ecosystem becomes contestable.\footnote{%
For more on digital inclusion, linked to both access and use of online spaces, see \citet{candidatu_leurs_ponzanesi2019}.} 
The ambivalence shown in this respect by both platforms and states ultimately disguises deeper structural processes of division at play among groups based on de facto inequalities \citep{nakamura2013}.

In the end, the effect of \isi{hate speech} primarily reflects attempts to shape social consensus on the acceptability of insulting people of different ethnicity or faith \citep[25]{kreissel_et_al2018}. This point accords well with speculative findings that online abuse on mainstream platforms that is serious enough to prompt internal criminal law action is low \citep[5]{vidgen_margetts_harris2019}.\footnote{%
The findings remain speculative because platforms do not share the amounts of content hosted.} 
The prediction of the collective harms at a societal dimension are not easy to quantify either in sociocultural terms. What is, however, already evident is that \isi{hate speech} of any nature contributes to the suppression of pluralism and inclusivity. In other words, \isi{hate speech} exposes our societies to a regression of the principles of equality and non-discrimination, including in relation to the protection of \isi{minorities} and freedom of religion in ways that are threatening to much broader groups than their immediate targets.

% \section*{Abbreviations}
% \begin{tabularx}{.5\textwidth}{@{}lQ@{}}
% CAT & Critical Discourse Analysis\is{Critical Discourse Analysis (CAT)} \\
% \end{tabularx}%
% \begin{tabularx}{.5\textwidth}{@{}lQ@{}}
% CDA & \isi{Critical Discourse Analysis} \\
% \end{tabularx}

% \section*{Acknowledgements}
% % \citet{Nordhoff2018} is useful for compiling bibliographies
% %\section*{Contributions}
% %John Doe contributed to conceptualization, methodology, and validation. 
% %Jane Doe contributed to writing of the original draft, review, and editing.

\sloppy
\printbibliography[heading=subbibliography,notkeyword=this]
\end{document}
